\section{循環する積分}\label{節:部分積分の循環}

指数関数と三角関数の積では,整式のように微分を繰り返して $0$ になる因子がない.では,部分積分を繰り返すとどうなるだろうか.一般的な形は,\,$a\neq0$ として\footnote{$a=0$ なら $\int\sin bx\dd x$ であり,部分積分は要らない.}
\[
 \int\e^{ax}\sin bx\dd x,\qquad
 \int\e^{ax}\cos bx\dd x
\]
で与えられる.正弦と余弦は微分によって互いに移り変わるので,2回の部分積分で,もとの積分が再び現れる.そこで計算を失敗と決めず,もとの積分を未知数として含む等式を作り,それを方程式として解くことを目標とする.

積分定数に惑わされずに等式を整理するため,まず
\[
 I(x)=\int_0^x\e^{at}\sin bt\dd t,\qquad
 J(x)=\int_0^x\e^{at}\cos bt\dd t
\]
と定める.\,$I,J$ は $x$ の関数であり,基本定理から $I'(x)=\e^{ax}\sin bx$ である.したがって $I(x)$ が求まれば,それに任意定数を加えて不定積分の答えにできる.

$I(x)$ で $\sin bt$ を微分し,\,$\e^{at}$ を積分すると,
\[
\begin{aligned}
 I(x)
 &=\left[\frac1a\e^{at}\sin bt\right]_0^x-\frac ba\int_0^x\e^{at}\cos bt\dd t\\
 &=\frac1a\e^{ax}\sin bx-\frac baJ(x)
\end{aligned}
\]
となる.続いて $J(x)$ でも $\e^{at}$ を積分する側に選ぶ.\,$\cos bt$ の微分は $-b\sin bt$ だから,
\[
\begin{aligned}
 J(x)
 &=\left[\frac1a\e^{at}\cos bt\right]_0^x+\frac ba\int_0^x\e^{at}\sin bt\dd t\\
 &=\frac1a\left(\e^{ax}\cos bx-1\right)+\frac baI(x)
\end{aligned}
\]
となる.ここで $I(x)$ が戻ってきた.第2式を第1式へ代入すると
\[
 I(x)
 =\frac1a\e^{ax}\sin bx-\frac b{a^2}\left(\e^{ax}\cos bx-1\right)-\frac{b^2}{a^2}I(x)
\]
となる.右辺の $I(x)$ を左辺へ移項し,両辺に $a^2$ を掛ければ
\[
 (a^2+b^2)I(x)=\e^{ax}(a\sin bx-b\cos bx)+b
\]
を得る.\,$a^2+b^2>0$ なので両辺をこれで割ることができ,
\[
 I(x)=\frac{\e^{ax}(a\sin bx-b\cos bx)}{a^2+b^2}+\frac{b}{a^2+b^2}
\]
となる.最後の定数を任意定数にまとめれば,不定積分が求まったのだ.

\begin{bookpoint}
\textbf{POINT}\par
指数関数を2回とも積分する側に選び,もとの積分 $I$ が戻った等式を $I$ についての方程式として解く.
\end{bookpoint}

\begin{example}{指数関数と正弦の積}
$\int\e^x\sin x\dd x$ を求めよ.
\end{example}

\paragraphbf{方針}
$I(x)=\int_0^x\e^t\sin t\dd t$,\,$J(x)=\int_0^x\e^t\cos t\dd t$ と置き,\,$\e^t$ を積分する部分積分を2回行って $I(x)$ についての方程式を作る.

\paragraphbf{解答}
$\frac12\e^x(\sin x-\cos x)+C$

\paragraphbf{解法}
$I(x)$ で $\sin t$ を微分し,\,$\e^t$ を積分すると,
\[
\begin{aligned}
 I(x)
 &=[\e^t\sin t]_0^x-\int_0^x\e^t\cos t\dd t\\
 &=\e^x\sin x-J(x)
\end{aligned}
\]
となる.\,$J(x)$ でも $\e^t$ を積分する側に選ぶと,
\[
\begin{aligned}
 J(x)
 &=[\e^t\cos t]_0^x-\int_0^x\e^t(-\sin t)\dd t\\
 &=\e^x\cos x-1+I(x)
\end{aligned}
\]
である.第2式を第1式へ代入すると
\[
 I(x)
 =\e^x\sin x-\{\e^x\cos x-1+I(x)\}
\]
となる.\,$I(x)$ を移項して
\[
 2I(x)=\e^x(\sin x-\cos x)+1
\]
を得る.よって
\[
 I(x)=\frac12\e^x(\sin x-\cos x)+\frac12
\]
である.最後の $\frac12$ も任意定数へまとめれば,
\[
 \int\e^x\sin x\dd x
 =\frac12\e^x(\sin x-\cos x)+C
\]
となる.これは上の一般形で $a=b=1$ とした場合に一致する.微分による確認は
\[
\begin{aligned}
 &\frac{\dd}{\dd x}\left\{\frac12\e^x(\sin x-\cos x)\right\}\\
 &\qquad=\frac12\e^x(\sin x-\cos x)
 +\frac12\e^x(\cos x+\sin x)\\
 &\qquad=\e^x\sin x
\end{aligned}
\]
である.

\paragraphbf{同じ積分が戻っても解けない場合}
ここで役立ったのは,「同じ積分が戻る」という事実だけではない.右辺へ $-I(x)$ として戻ったので,移項すると $2I(x)$ が得られたことが肝である.一般形でも,戻ってきた $I(x)$ の係数は $-\frac{b^2}{a^2}$ であり,移項した後の係数 $1+\frac{b^2}{a^2}$ は $0$ にならなかった.

一方,\,$J(x)$ の部分積分で,\,$\e^t$ を微分する側に選び直したとしよう.すると1回目の計算を逆向きにたどることになり,\,$I(x)=I(x)$ という等式に戻ってしまう.これは正しい等式だが,\,$I(x)$ についての新しい情報はない.循環した後には,何倍の $I(x)$ が残ったのかを確かめるとよい.

\paragraphbf{不定積分のまま計算する場合}
不定積分のまま同じ計算を進めることもできる.その場合は,途中の等式が任意定数の違いを許していることを忘れず,最後に1つの $C$ にまとめる.定積分として $I,J$ を定めたのは,等式の両辺を同じ関数として固定し,移項の意味をはっきりさせるためである.
